\chapter[1977 --- Guillemin et al.]{1977: Roger Guillemin, Andrew V. Schally, Rosalyn Yalow}
\label{ch:1977_Guillemin}

\section*{Main Contribution}

\textit{Discovered the hypothalamic hormones that control the pituitary gland, revealing how the brain regulates growth, reproduction, and metabolism.}

\section{Official Citation}

for their discoveries concerning the peptide hormone production of the brain

\section[Roger Guillemin and Andrew V.]{Roger Guillemin and Andrew V. Schally: Discoveries Concerning the Peptide Hormone Production of the Brain}

\subsection{Background and Historical Context}

The 1977 Nobel Prize in Physiology or Medicine was divided: one half was awarded jointly to Roger Guillemin and Andrew V. Schally ``for their discoveries concerning the peptide hormone production of the brain,'' and the other half was awarded to Rosalyn Yalow ``for her development of the radioimmunoassay of peptide hormones.'' Together, their work revolutionized our understanding of how the brain controls the endocrine system and created powerful new tools for measuring hormones in the blood.

For much of the twentieth century, the relationship between the brain and the endocrine system was poorly understood. It was known that the pituitary gland, a small organ at the base of the brain, produced hormones that regulated growth, reproduction, metabolism, and other vital functions. It was also known that the hypothalamus, a region of the brain just above the pituitary, somehow controlled the pituitary's activity. However, the mechanism by which the hypothalamus communicated with the pituitary---and the nature of the hypothalamic ``releasing hormones'' that regulated pituitary function---remained unknown.

Roger Guillemin was born in Dijon, France, in 1924, and trained in medicine at the University of Lyon. He emigrated to the United States in 1949 and joined the laboratory of Hans Selye in Montreal before moving to the Baylor College of Medicine in Houston, Texas. Andrew V. Schally was born in Wilno, Poland (now Vilnius, Lithuania), in 1926, and trained in medicine at the University of London. He also emigrated to the United States, joining the Baylor College of Medicine in 1957, where he and Guillemin were colleagues for a period before Schally moved to the Veterans Administration Hospital in New Orleans.

The race to identify the hypothalamic releasing hormones was one of the most intensely competitive episodes in the history of endocrinology. Guillemin and Schally, both trained in the French tradition of biochemistry, were the principal rivals in this contest, and their competition drove both men to notable lengths of effort and ingenuity.

\subsection{The Discovery/Contribution}

The identification of the hypothalamic releasing hormones was a monumental biochemical challenge. The hypothalamus produces these hormones in vanishingly small quantities---on the order of nanograms or even picograms per hypothalamus---and isolating and characterizing them required processing enormous numbers of animal hypothalami. Guillemin and Schally each organized massive extraction projects, collecting hundreds of thousands of sheep and pig hypothalami from slaughterhouses to obtain enough material for biochemical analysis.

Schally and his team at the New Orleans VA Hospital were the first to succeed. In 1969, Schally's group isolated and characterized thyrotropin-releasing hormone (TRH), a small peptide of only three amino acids (pyroGlu-His-ProNH$_2$) that stimulates the release of thyroid-stimulating hormone (TSH) from the anterior pituitary. The determination of TRH's structure was a tour de force of biochemical methodology, requiring the purification of approximately one million porcine hypothalami to obtain one milligram of the active peptide.

Shortly after the discovery of TRH, Schally's group went on to isolate and characterize luteinizing hormone-releasing hormone (LHRH, later renamed GnRH for gonadotropin-releasing hormone) in 1971. LHRH is a decapeptide that controls the release of luteinizing hormone (LH) and follicle-stimulating hormone (FSH) from the pituitary, thereby regulating reproductive function. The structure of LHRH---pGlu-His-Trp-Ser-Tyr-Gly-Leu-Arg-Pro-GlyNH$_2$---was determined by Edman degradation and confirmed by total chemical synthesis.

Guillemin's group, working at the Salk Institute in La Jolla, made parallel discoveries. Guillemin isolated corticotropin-releasing hormone (CRH), which controls the release of adrenocorticotropic hormone (ACTH) from the pituitary and thereby regulates the stress response. Guillemin also contributed to the characterization of growth hormone-releasing hormone (GHRH), which stimulates the release of growth hormone from the anterior pituitary.

The identification of these hypothalamic peptides established a new principle in endocrinology: that the brain produces hormones that regulate the pituitary gland, and that these hormones are small peptides that act as chemical messengers between the nervous system and the endocrine system. This discovery---the concept of neuroendocrine regulation---fundamentally changed our understanding of how the body maintains hormonal homeostasis. It also opened the door to the development of new diagnostic and therapeutic tools for disorders of the hypothalamic-pituitary axis.

\section[Rosalyn Yalow: Development of the]{Rosalyn Yalow: Development of the Radioimmunoassay of Peptide Hormones}

\subsection{Background and Historical Context}

Rosalyn Sussman Yalow was born in New York City in 1921 and trained in physics at the University of Illinois, earning her Ph.D. in 1945. She joined the Veterans Administration Hospital in the Bronx in 1950, where she worked with Solomon Berson on the study of insulin metabolism in diabetic patients. Yalow's background in physics gave her a unique perspective on biological measurement, and she brought the quantitative rigor of physics to the study of hormone biology.

\subsection{The Discovery/Contribution}

The radioimmunoassay (RIA) technique that Yalow developed with Berson was one of the major methodological innovations in the history of medicine. The technique was born out of a seemingly simple observation: when insulin was injected into patients, some of the insulin was bound by antibodies in the blood. Berson and Yalow realized that this antibody binding could be exploited to measure the concentration of insulin in blood samples with notable sensitivity and specificity.

The principle of the radioimmunoassay is elegant in its simplicity. A fixed amount of radioactively labeled hormone (the ``tracer'') is mixed with a fixed amount of antibody specific for that hormone, along with varying amounts of unlabeled hormone (the ``standard'' or the unknown sample). The labeled and unlabeled hormones compete for binding to the antibody. The more unlabeled hormone present, the less labeled hormone is bound. By measuring the fraction of labeled hormone bound to the antibody, the concentration of unlabeled hormone in the sample can be determined with great precision.

Yalow and Berson published their first description of the radioimmunoassay for insulin in 1959, and the technique was quickly recognized as a significant advance. Before the development of RIA, it was impossible to measure the concentration of insulin in blood with any accuracy, because insulin is present in such small quantities. The radioimmunoassay was sensitive enough to measure hormone concentrations in the picogram-per-milliliter range---a level of sensitivity that was notable in clinical chemistry.

Yalow's group went on to develop radioimmunoassays for a wide range of hormones and other biologically active substances, including growth hormone, parathyroid hormone, ACTH, angiotensin, cyclic AMP, and many others. The technique was also adapted for the measurement of drugs, vitamins, and viral antigens, including hepatitis B surface antigen---a test that became crucial for screening blood donations.

Yalow was awarded the Nobel Prize in 1977 for her development of the radioimmunoassay. She was the first woman to receive the Nobel Prize in Physiology or Medicine without sharing it with a co-laureate (though she shared the prize year with Guillemin and Schally, her prize was specifically for the RIA technique). She was also the second woman to receive the Nobel Prize in Physiology or Medicine, after Gerty Cori in 1947.

Berson, Yalow's collaborator and close scientific partner, had died of cancer in 1972, five years before the prize was awarded. Nobel Prizes are not given posthumously, so Berson could not share in the honor, but Yalow always emphasized that the development of the radioimmunoassay was a collaborative achievement.

\subsection{Impact on Modern Medicine}

The impact of the radioimmunoassay on modern medicine is difficult to overstate. The technique transformed clinical endocrinology by making it possible to measure hormones in the blood with accuracy and precision for the first time. Before RIA, clinicians had to rely on indirect and often unreliable methods to assess hormone levels---such as measuring the biological effects of hormones (e.g., blood sugar levels for insulin, or growth rates for growth hormone) rather than measuring the hormones themselves. The radioimmunoassay made it possible to diagnose endocrine disorders---such as diabetes, thyroid disease, growth hormone deficiency, and pituitary tumors---by directly measuring hormone concentrations in the blood.

The technique also had a significant impact on basic research. By enabling the precise measurement of hormone levels in blood and tissues, RIA allowed researchers to elucidate the regulation of hormone secretion, the pharmacokinetics of hormone administration, and the mechanisms of endocrine disease. The development of radioimmunoassays for neurotransmitters and neuropeptides also facilitated the study of brain chemistry and the development of psychopharmacology.

Beyond endocrinology, the radioimmunoassay found applications in virtually every area of medicine. RIA-based tests were developed for the measurement of therapeutic drugs (e.g., digoxin, theophylline, and lithium), enabling more precise dosing and reducing the risk of toxicity. RIA was used to detect viral antigens (e.g., hepatitis B surface antigen) and antibodies (e.g., anti-HIV antibodies), forming the basis of blood screening tests that have saved millions of lives. RIA was also used in forensic medicine, environmental monitoring, and food safety testing.

The development of the radioimmunoassay also paved the way for the later development of non-radioactive immunoassay techniques, including enzyme-linked immunosorbent assays (ELISA) and chemiluminescent immunoassays, which have largely replaced RIA in clinical laboratories. However, the fundamental principles of competitive binding assays that Yalow and Berson established remain the basis of all modern immunoassay techniques. Yalow's contribution to medicine is thus not limited to a single technique; it encompasses the entire field of immunoassay and the vast array of diagnostic tests that have been developed using its principles.

\section{Legacy}

The work of Guillemin, Schally, and Yalow transformed endocrinology from a largely observational discipline into a precise, molecular science. Guillemin and Schally's identification of the hypothalamic releasing hormones revealed the intricate chemical communication between the brain and the endocrine system, providing the foundation for modern neuroendocrinology and leading to new treatments for disorders of the hypothalamic-pituitary axis. Synthetic analogs of GnRH are now widely used in the treatment of prostate cancer, endometriosis, and infertility, and somatostatin analogs (derived from the discovery of the hypothalamic inhibitory hormones) are used in the treatment of acromegaly and certain neuroendocrine tumors.

Yalow's radioimmunoassay, meanwhile, became one of the most widely used techniques in the history of medicine. It was applied to thousands of different substances and played a central role in the development of clinical chemistry, pharmacology, and diagnostic medicine. Yalow herself continued to work at the VA Hospital in the Bronx until her retirement in 1992, and she became a prominent advocate for women in science and a vocal critic of pseudoscience.

Together, the work of the three 1977 laureates demonstrated the power of biochemical and biophysical approaches to solve fundamental problems in medicine. Their discoveries about the peptide hormones of the brain and the techniques for measuring them have saved countless lives and continue to shape our understanding of the endocrine system and its role in health and disease.


\section{Modern Developments}

Guillemin, Schally, and Yalow's work on hypothalamic hormones has led to modern endocrinology. Understanding of GnRH has enabled treatments for prostate cancer (leuprolide), endometriosis, and precocious puberty. GH-releasing hormone analogs treat growth hormone deficiency. Understanding of somatostatin has led to octreotide for acromegaly and neuroendocrine tumors. Peptide-based drugs, including GLP-1 receptor agonists for diabetes and obesity, build on their discoveries.
